%======================================================================
\section{Local physics}\label{sec:physics}
%======================================================================

Theorem~B concerns an abstract channel. Here we record what it implies for physical systems that could produce the channel, and what
autonomous dynamics can and cannot produce in the long run. We call a channel \emph{cheap} if for all $n$ and $\epsilon$ it has causal
services for $n$ rounds with error at most $\epsilon$, $P=C=0$ and memory polynomial in $\log(n/\epsilon)$, and \emph{expensive} if for some
$A,c>0$ every service for $n$ rounds with error at most $1/16$ and $P+C\le A\log(n+1)$ needs memory at least $n^c$ for all large $n$. Neither notion
constrains the exchange.

\paragraph{Local reservoirs.} A \emph{local reservoir} consists of a finite box of the lattice $\Z^\nu$, with a system of dimension at most $q$ at
each site, the probe $\mathbb C^d$ attached to the origin, and a Hamiltonian $H(t)=\sum_Xh_X(t)$, possibly time dependent, whose terms are supported
in sets $X$ of diameter at most $R_0$ and satisfy $\|h_X(t)\|\le J_0$. The reservoir starts in its tracial state, the probe is coupled to it for a
time $T$, and $\Phi_T(X)=(\mathrm{id}\otimes\tau)[W_T(X\otimes I)W_T^*]$ is the resulting channel on the probe, where $W_T$ is the time-ordered
evolution and $\tau$ the normalized trace of the reservoir. We call $\nu,q,R_0,J_0$ and $d$ the \emph{parameters} of the reservoir; the size of the
box is arbitrary.

\begin{proposition}[Local reservoirs are cheap at every finite time]\label{prop:light-cone}
There is $K$, depending only on the parameters, such that for every local reservoir, every $T\ge0$, every $n\ge1$ and every $0<\epsilon\le1$, the
channel $\Phi_T$ has a causal service for $n$ rounds with error at most $\epsilon$, $P=C=0$, $J=(n-1)Q$ and
\[
Q\ \le\ K\,\bigl[1+T+\log(n/\epsilon)\bigr]^\nu .
\]
\end{proposition}

\begin{proof}
For $r\ge0$ let $\Phi_T^{(r)}$ be the channel produced by the terms $h_X$ with $X$ inside the ball $B_r$ of radius $r$ about the origin. The
Lieb--Robinson bound~\cite{LiebRobinson1972,BHV2006,NSY2019} gives constants $C_1,v,\mu>0$, depending only on the parameters, such that the Heisenberg
evolutions of an operator $A$ at the origin under the full and the truncated dynamics differ in norm by at most $C_1\|A\|e^{vT-\mu r}$. The bound
does not depend on the dimension of the site at the origin, so it holds for operators on the probe tensored with a reference system that the
dynamics does not touch. Hence the Heisenberg-picture maps of $\Phi_T$ and $\Phi_T^{(r)}$ differ in completely bounded norm by at most
$C_1e^{vT-\mu r}$, and by duality $\ddia(\Phi_T,\Phi_T^{(r)})\le C_1e^{vT-\mu r}$. The truncated evolution acts only on the probe and the sites
of $B_r$, and the other sites start in their tracial state, so $\Phi_T^{(r)}$ has a finite tracial factorization with bath dimension
$m\le q^{(2r+1)^\nu}$.

Choose $r=\lceil(vT+\ln(2C_1n/\epsilon))/\mu\rceil$, so that $\ddia(\Phi_T,\Phi_T^{(r)})\le\epsilon/(2n)$, and let $Q=\lceil\log m+\log(2n/\epsilon)\rceil$.
The exchange construction in the proof of Proposition~\ref{prop:sofic-upper}, with $U_f$ replaced by the unitary of this factorization and
$\vartheta=0$, gives a service of $\Phi_T^{(r)}$ in which every round is within $\epsilon/(2n)$ of $\Phi_T^{(r)}$, hence within $\epsilon/n$ of
$\Phi_T$; a hybrid argument bounds the error by $\epsilon$. Finally $Q\le(2r+1)^\nu\log q+\log(2n/\epsilon)+1\le K[1+T+\log(n/\epsilon)]^\nu$.
\end{proof}

With Theorem~B this says that the channel of the configuration lamps cannot be produced quickly by local physics.

\begin{corollary}[Preparation time]\label{cor:preparation}
Let $\Phi$ be the channel of Theorem~\ref{thm:main-channel}, and fix the parameters of a local reservoir with $d=873$. There are $c,\delta_0>0$ such
that, whenever such a reservoir produces at time $T$ a channel within $\delta\le\delta_0$ of $\Phi$ in diamond distance,
\[
T\ \ge\ c\,\delta^{-1/\nu}\bigl(\log(1/\delta)\bigr)^{-(2\kappa+1)/\nu},\qquad 2\kappa+1<26.49 .
\]
\end{corollary}

\begin{proof}
Let $n=\lfloor1/(32\delta)\rfloor$. Proposition~\ref{prop:light-cone} gives a service of $\Phi_T$ with error at most $1/32$, zero purity and
$Q\le K[1+T+\log(32n)]^\nu$. Replacing the $n$ uses of $\Phi_T$ by uses of $\Phi$ one at a time changes the observer's final state by at most
$n\delta\le1/32$, so the same device serves $\Phi$ with error at most $1/16$ and zero purity. For $\delta_0$ small, $n$ is large, and
Theorem~\ref{thm:main-channel}(a) with $A_0=0$ gives $Q\ge c'n/(\log n)^{2\kappa+1}$. Hence
$1+T+\log(32n)\ge\bigl(c'n/(K(\log n)^{2\kappa+1})\bigr)^{1/\nu}$, and since $n\ge1/(64\delta)$ this gives the claim for $\delta_0$ small.
\end{proof}

The maximally mixed start of the reservoir is essential: $\Phi$ has a Stinespring dilation with a pure environment of dimension $r_*$, so
from a pure state a bounded number of sites produce $\Phi$ exactly in bounded time. The same argument applies to a circuit of depth $T$ of gates of
bounded range, whose light cone is exact, and so bounds the depth of any local
circuit that prepares $\Phi$. The corollary does not exclude that local dynamics produce $\Phi$ in the limit of long times; it says how slowly.

\paragraph{Autonomous dynamics.} Let $\mathcal B$ be a von Neumann algebra with a faithful normal tracial state $\tau_{\mathcal B}$, put
$\mathcal M=M_d\otimes\mathcal B$ with the trace $\tau=\mathrm{tr}_d\otimes\tau_{\mathcal B}$, where $\mathrm{tr}_d$ is the normalized trace, and let
$\iota(a)=a\otimes1$ and $E=\mathrm{id}\otimes\tau_{\mathcal B}$. An \emph{autonomous dynamics} is a trace-preserving automorphism $\alpha$ of $\mathcal M$,
and the channel after $t\in\N$ steps is $\Phi_t=E\circ\alpha^t\circ\iota$; for $\alpha=\mathrm{Ad}\,U$ with $\mathcal B=M_m$ this is the tracial
factorization $\Phi_t(X)=(\mathrm{id}\otimes\tau_m)[U^t(X\otimes I)U^{-t}]$. Let $\mathcal N$ be the fixed-point algebra of $\alpha$ and
$E_{\mathcal N}$ the $\tau$-preserving conditional expectation onto it.

The next proposition combines the mean ergodic theorem with the observation of Junge, Le Merdy and Xu~\cite[Remark~10.3]{JungeLeMerdyXu2006}
and of Haagerup and Musat~\cite[Remark~5.2]{HaagerupMusat2011} that compressions $E\circ E_{\mathcal N}\circ\iota$ of trace-preserving conditional
expectations are self-adjoint and positive; both of its clauses are standard.

\begin{proposition}[{Autonomous limits are self-adjoint; the mean ergodic theorem and~\cite[Remark~10.3]{JungeLeMerdyXu2006},
\cite[Remark~5.2]{HaagerupMusat2011}}]\label{prop:autonomous}
The averages $\frac1T\sum_{t=0}^{T-1}\Phi_t$ converge in diamond norm to $\overline\Phi=E\circ E_{\mathcal N}\circ\iota$, and if $\Phi_t$ converges, its
limit is $\overline\Phi$. The channel $\overline\Phi$ is self-adjoint and positive semidefinite for the Hilbert--Schmidt inner product on $M_d$. In
particular, a channel $\Psi$ with $\Psi^*\ne\Psi$ is not the long-time limit of any autonomous dynamics with a tracial bath, local or not.
\end{proposition}

\begin{proof}
The automorphism extends to a unitary $V$ on $L^2(\mathcal M,\tau)$. By the mean ergodic theorem~\cite{Krengel1985}, $\frac1T\sum_{t<T}V^t$ converges
strongly to the orthogonal projection $P$ onto the invariant vectors. For $x\in\mathcal M$ the averages of $x$ lie in the ball of radius $\|x\|$ of
$\mathcal M$, which is closed in $L^2$, so $Px\in\mathcal M$, and $Px$ is invariant; hence $P(\mathcal M)\subseteq\mathcal N$. Since
$\mathcal N$ consists of invariant vectors and $\mathcal M$ is dense in $L^2(\mathcal M)$, $P$ is the orthogonal projection onto
$L^2(\mathcal N)$, whose restriction to $\mathcal M$ is $E_{\mathcal N}$. The map $\iota$ is an isometry from $(M_d,\mathrm{tr}_d)$ into
$L^2(\mathcal M,\tau)$ with adjoint $E$, so $\Phi_t=\iota^*V^t\iota$, and the averages converge to $\iota^*P\iota=\overline\Phi$ on every $a\in M_d$,
hence in diamond norm, since $M_d$ is finite dimensional. An ordinary limit equals the limit of the averages. Finally
$\mathrm{tr}_d(b^*\overline\Phi(a))=\langle P\iota(b),P\iota(a)\rangle$ is Hermitian in $(a,b)$ and nonnegative for $a=b$.
\end{proof}

Self-adjointness and positivity are not the only constraints. The limit $\overline\Phi$ is a compression of a conditional expectation, and this ties
the decay of a coherence to the return probability of a state.

\begin{proposition}[A return inequality]\label{prop:return}
Let $T=E\circ E_{\mathcal N}\circ\iota$ be as in Proposition~\ref{prop:autonomous}, on $M_D$. If $T(e_{00})=e_{00}$ and $T(e_{0j})=a\,e_{0j}$ for some
$j\ne0$ and $0\le a\le1$, then
\[
\langle j|T(e_{jj})|j\rangle\ \ge\ a^2+\frac{(1-a)^2}{D-1}.
\]
\end{proposition}

\begin{proof}
Write $P=E_{\mathcal N}$, the orthogonal projection onto $L^2(\mathcal N)$, and identify $x\in M_D$ with $\iota(x)$; then
$\mathrm{tr}_D(x^*T(y))=\langle x,Py\rangle$ with $\langle x,y\rangle=\tau(x^*y)$. Put $p_0=e_{00}$, $p_j=e_{jj}$ and $x=e_{0j}$. Since
$\|Pp_0\|_2^2=\langle p_0,T(p_0)\rangle=\|p_0\|_2^2$, we have $Pp_0=p_0$, so $p_0\in\mathcal N$. Let $y=P(x)$, $X=P(p_j)$ and $B=y^*y$. As $P$ is
$\mathcal N$-bimodular, $p_0y=P(p_0x)=y$, $yp_0=P(xp_0)=0$ and $p_0X=P(p_0p_j)=0$; so $yy^*$ is supported in $p_0$, and $B$ and $X$ in $1-p_0$. The
Schwarz inequality for conditional expectations gives $B\le P(x^*x)=X$. Moreover $\tau(X)=\tau(p_j)=1/D$ and $\tau(B)=\|Px\|_2^2=\langle x,T(x)\rangle=a/D$,
so $R=X-B\ge0$ has $\tau(R)=(1-a)/D$. By the Cauchy--Schwarz inequality, $\tau(B^2)=\tau((yy^*)^2)\ge\tau(yy^*)^2/\tau(p_0)=a^2/D$ and
$\tau(R^2)\ge\tau(R)^2/\tau(1-p_0)=(1-a)^2/(D(D-1))$, and $\tau(BR)\ge0$. Hence
\[
\frac1D\,\langle j|T(e_{jj})|j\rangle=\langle p_j,Pp_j\rangle=\tau(X^2)=\tau(B^2)+2\tau(BR)+\tau(R^2)\ \ge\ \frac{a^2}D+\frac{(1-a)^2}{D(D-1)} .\qedhere
\]
\end{proof}

\begin{remark}[Self-adjoint and positive is not enough]\label{rem:qutrit}
Haagerup and Musat give self-adjoint positive maps that are not such compressions because they are not
factorizable~\cite[Theorem~5.4]{HaagerupMusat2011}. The following channel is an exact finite factor and is still excluded. Let $a=\sqrt3/2$ and $b=\frac12$, and let $\Lambda$ on $M_3$ have the Kraus operators $\mathrm{diag}(1,a,a)$, $b\,e_{12}$ and $-b\,e_{21}$. With
$\omega=e^{2\pi i/3}$ and $W_k=1\oplus\bigl(\begin{smallmatrix}a&b\omega^k\\-b\omega^{-k}&a\end{smallmatrix}\bigr)$ one has
$\Lambda=\frac13\sum_{k=0}^2\mathrm{Ad}\,W_k$, so $\Lambda$ has a finite tracial factorization with bath dimension three. It is self-adjoint, and positive
for the Hilbert--Schmidt inner product, with eigenvalues between $\frac12$ and $1$. But $\Lambda(e_{00})=e_{00}$, $\Lambda(e_{01})=a\,e_{01}$ and
$\langle1|\Lambda(e_{11})|1\rangle=a^2$, which violates Proposition~\ref{prop:return}. So $\Lambda$ is not the long-time limit of any autonomous dynamics
with a tracial bath.
\end{remark}

Conversely, every channel of the form $\Psi\Psi^*$ with $\Psi$ in the tracial closure is a compression of a conditional
expectation (apply~\cite[Theorem~5.9]{HaagerupMusat2011} to $\Psi^*$, which is again factorizable), and in fact an autonomous long-time limit, of a dynamics that the construction does not make
local; the construction adapts Kümmerer's tensor dilations with a Bernoulli shift~\cite{Kummerer1985}.

\begin{proposition}[Autonomy without locality]\label{prop:nonlocal}
Let $\mathcal N_0$ be a von Neumann algebra with a faithful normal tracial state $\tau_0$, let $U$ be a unitary in $M_d\otimes\mathcal N_0$, and put
$\Psi(X)=(\mathrm{id}\otimes\tau_0)[U(X\otimes1)U^*]$. Let $\mathcal B=\overline\bigotimes_{k\in\Z}(\mathcal N_0,\tau_0)$, with $U$ acting on the factor
$k=0$, and let $\beta$ be the Bernoulli shift of $\mathcal B$. Then the autonomous dynamics $\alpha=\mathrm{Ad}\,U\circ(\mathrm{id}\otimes\beta)\circ\mathrm{Ad}\,U^*$
has $\Phi_t\to\Psi\Psi^*$, where $\Psi^*(a)=(\mathrm{id}\otimes\tau_0)[U^*(a\otimes1)U]$ is the Hilbert--Schmidt adjoint of $\Psi$.
\end{proposition}

\begin{proof}
We have $\alpha^t=\mathrm{Ad}\,U\circ(\mathrm{id}\otimes\beta^t)\circ\mathrm{Ad}\,U^*$. The shift is mixing: $\tau_{\mathcal B}(w\,\beta^t(z))\to\tau_{\mathcal B}(w)\tau_{\mathcal B}(z)$ for
$w,z\in\mathcal B$, since equality holds as soon as $w$ and $\beta^t(z)$ are elementary tensors with disjoint supports, and such tensors of norm
at most $\|w\|$ and $\|z\|$ approximate $w$ and $z$ in $2$-norm (Kaplansky's density theorem). Let $a\in M_d$ and $y=U^*(a\otimes1)U$, and write
$y=\sum_{i,j}e_{ij}\otimes y_{ij}$. For every $x=\sum_{k,l}e_{kl}\otimes x_{kl}$ in $M_d\otimes\mathcal B$,
\[
\tau\bigl(x\,(\mathrm{id}\otimes\beta^t)(y)\bigr)=\frac1d\sum_{k,l}\tau_{\mathcal B}\bigl(x_{kl}\beta^t(y_{lk})\bigr)\ \longrightarrow\ \tau\bigl(x\,(E(y)\otimes1)\bigr).
\]
Taking $x=U^*(e_{ji}\otimes1)U$ shows that every matrix entry of $\Phi_t(a)=E[U(\mathrm{id}\otimes\beta^t)(y)U^*]$ converges to the corresponding entry
of $E[U(E(y)\otimes1)U^*]=\Psi(\Psi^*(a))$. Convergence on a basis of $M_d$ gives convergence in diamond norm.
\end{proof}

By a theorem of Haagerup and Musat~\cite[Theorem~3.6]{HaagerupMusat2015}, every channel in $\cK_d$ has the form $\Psi$ of
Proposition~\ref{prop:nonlocal} for some $\mathcal N_0$, and the channel of Theorem~B has it with $\mathcal N_0=L(\Gamma)$.

Nothing in Proposition~\ref{prop:nonlocal} makes the dynamics local: $U$ is an arbitrary unitary of $M_d\otimes\mathcal N_0$, and we do not know
whether a dynamics of this kind can be given by a local rule.

\begin{remark}[Self-adjoint expensive channels]\label{rem:doubling}
Self-adjointness is no obstacle to expense. Let $z$ be the unitary of order three in $L(\Z_3)$, so that $\tau(z)=\tau(z^2)=0$, let $U$ be the gadget
unitary of the channel $\Phi$ of Theorem~B, and let
$V=\bigl(\begin{smallmatrix}0&U\otimes z\\U^*\otimes z^*&0\end{smallmatrix}\bigr)$, a self-adjoint unitary over $\mathbb C^2\otimes\mathbb C^{873}$ and
$L(\Gamma\times\Z_3)$. The channel it defines on $M_{1746}$ is
$\Theta\bigl(\begin{smallmatrix}x&y\\u&w\end{smallmatrix}\bigr)=\mathrm{diag}\bigl(\Phi(w),\Phi^*(x)\bigr)$, which is self-adjoint. It keeps the lower bounds of
Theorem~B. For part~(a), a service of $\Theta$ becomes a service of $\Phi$ with one more memory qubit and one more bit of purity, at the same
exchange, by feeding each input together with a retained tag qubit that starts in $|1\rangle$ and is flipped after every round, since $\Theta$ maps
$|1\rangle\langle1|\otimes\rho$ to $|0\rangle\langle0|\otimes\Phi(\rho)$. For part~(b), $\Theta$ has a flat probe of rank $2r_*$, with anchors in both blocks,
so Proposition~\ref{prop:purity-budget} gives $P+C<3Q+6$ at its own rank rate $J\le\frac n2\log(2r_*)$; the tag transfer and part~(a) then give
$Q\ge c\sqrt n/(\log n)^\kappa$. But $\Theta$ maps $\mathrm{diag}(I,-I)$ to $-\mathrm{diag}(I,-I)$, so it is not
positive and not an autonomous limit, while $\Theta^2=\mathrm{diag}(\Phi\Phi^*,\Phi^*\Phi)$ is one, by Proposition~\ref{prop:nonlocal}. The next
proposition shows that $\Theta^2$ is cheap.
\end{remark}

Composing a gadget channel with its adjoint erases the multiplication of the group. In a triangle block the unitarity of the gadget forces the
coefficient of the product label to be the product of the coefficients, and this is what the one-round theorem reads. In $\Phi\Phi^*$ a triangle only
requires a permutation $\pi_x$ of the labels that carries $1$ to $x$ and $y$ to $xy$, and no relation between the permutations $\pi_x$ for
different $x$ is needed.

\begin{proposition}[Adjoint compositions of gadgets are exact finite factors]\label{prop:compositions}
Let $\Phi_G$ be a gadget channel with $r_*$ labels. Then $\Phi_G\Phi_G^*$ and $\Phi_G^*\Phi_G$ have finite tracial factorizations with baths of at
most $3r_*-2$ qubits. Consequently, for every $n$, they have exact causal services with $Q\le3r_*-2$, $P=C=0$ and $J=(n-1)Q$, and so does
$\Theta^2$ in Remark~\ref{rem:doubling}, with one more qubit.
\end{proposition}

\begin{proof}
Let $\mathcal F$ be the set of labels, and write $A_g=\sum_{[w_a]=g}c_aE_a$, so that $\Phi_G\Phi_G^*(X)=\sum_{g,h}A_gA_h^*XA_hA_g^*$. A product $A_gA_h^*$
is a sum of terms $c_a\overline{c_b}E_aE_b^*$ over slots $a,b$ in a common column, so it vanishes unless $g$ and $h$ label two slots of one column. In a
triangle block with labels $1,y$ in the first row and $x,z$ in the second, where $z=xy$, the columns carry the labels $\{1,x\}$ and $\{y,z\}$; so
$A_gA_h^*\ne0$ forces $hg^{-1}\in K=\{1\}\cup\{x^{\pm1}\}$, the union over the triangles. For $k\in K$ let $\pi_k$ be a permutation of $\mathcal F$ that
extends the partial map $g\mapsto kg$ (defined when $g,kg\in\mathcal F$), with $\pi_1=\mathrm{id}$ and $\pi_{k^{-1}}=\pi_k^{-1}$: for $k\ne k^{-1}$ extend
one of the two maps and invert it, and for an involution $k$ the partial map is an involution on its domain, which we extend by the identity.

The bath is $r_*+|K|-1\le3r_*-2$ qubits. On the first $r_*$, indexed by $\mathcal F$, let $V_g$ be the Pauli matrix $Z$ on the qubit $g$, so that
$\tau(V_gV_h)=\delta_{gh}$, and let $P_k$ permute these qubits by $\pi_k$, so that $P_kV_gP_k^*=V_{kg}$ whenever $g,kg\in\mathcal F$. On the other
qubits choose diagonal unitaries $\chi_k$ with $\chi_1=I$, $\chi_{k^{-1}}=\chi_k^*$ and $\tau(\chi_k^*\chi_l)=\delta_{kl}$: for an involution, $Z$ on one
qubit, and for a pair $k\ne k^{-1}$, $\mathrm{diag}(1,i,-1,-i)$ and its adjoint on two qubits, different pairs on different qubits. Put
$R_k=P_k\otimes\chi_k$, so that $R_{k^{-1}}=R_k^*$ and $R_kV_gR_k^*=V_{kg}$ for $g,kg\in\mathcal F$, and put $C_{g,h}=V_gR_{hg^{-1}}^*$. These are
orthonormal for the normalized trace: if $hg^{-1}\ne h'g'^{-1}$ the factor $\chi$ has trace zero, and otherwise
$\tau(C_{g,h}^*C_{g',h'})=\tau(R_kV_gV_{g'}R_k^*)=\delta_{gg'}$, which then forces $h=h'$. In the same way
$\tau(C_{g,h}C_{g',h'}^*)=\tau(V_gV_{g'})=\delta_{gg'}$ when $hg^{-1}=h'g'^{-1}$, and it vanishes otherwise.

Let $W=\sum A_gA_h^*\otimes C_{g,h}$, summed over the pairs with $hg^{-1}\in K$; the other products $A_gA_h^*$ vanish. On the block of a symbol
with label $g$, including the first block, where $g=1$, it is $V_g$. On a triangle block, with $H$ the Hadamard matrix and using
$R_xV_1R_x^*=V_x$ and $R_xV_yR_x^*=V_z$,
\[
W=\frac12\begin{pmatrix}V_1+V_y&(V_1-V_y)R_x^*\\ R_x(V_1-V_y)&V_x+V_z\end{pmatrix}
=\begin{pmatrix}I&0\\0&R_x\end{pmatrix}(H\otimes I)\begin{pmatrix}V_1&0\\0&V_y\end{pmatrix}(H\otimes I)\begin{pmatrix}I&0\\0&R_x^*\end{pmatrix},
\]
a self-adjoint unitary; triangles with coinciding labels are covered by the same computation. So $W$ is unitary, and by orthonormality of the
$C_{g,h}$, $(\mathrm{id}\otimes\tau)[W(X\otimes I)W^*]=\sum_{g,h}A_gA_h^*XA_hA_g^*=\Phi_G\Phi_G^*(X)$. For $\Phi_G^*\Phi_G$ use right translations
$g\mapsto gk$ with $k\in\{1\}\cup\{y^{\pm1}\}$, the operators $C_{g,h}=V_gR_{g^{-1}h}^*$, and rows instead of columns; on a triangle block the rows carry
$\{1,y\}$ and $\{x,z\}$, and the same computation applies with $x$ and $y$ exchanged. For $\Theta^2$, one further maximally mixed qubit dephases the
two blocks, and controlled on the block the two factorizations act on a common bath. Refreshing the whole bath in every round gives the services.
\end{proof}

These services refresh the whole bath in every round, so their exchange is far above the rank rate; whether $\Phi\Phi^*$ has services with $P=C=0$ and
memory polynomial in $\log(n/\epsilon)$ at its own rank rate we do not know.

So a self-running system with a tracial bath cannot reach expense through $\Phi\Phi^*$. It can through a lazy version of the doubling.

\paragraph{A lazy doubling.} The doubled channel $\Theta$ of Remark~\ref{rem:doubling} is expensive but not positive, and its square is cheap.
Averaging $\Theta$ with the identity restores positivity and keeps the expense. Put
\[
\Theta_+=\tfrac12(\mathrm{id}+\Theta),\qquad
\Theta_+\begin{pmatrix}x&y\\u&w\end{pmatrix}=\frac12\begin{pmatrix}x+\Phi(w)&y\\u&w+\Phi^*(x)\end{pmatrix}\qquad\text{on }M_{1746}.
\]
In Remark~\ref{rem:doubling} only $zz^*=1$ and $\tau(z^2)=0$ are used, so $z$ may also be the unitary $\mathrm{diag}(1,i)$ of a qubit with its
normalized trace. Since $\Theta$ is a unital channel that is self-adjoint for the Hilbert--Schmidt inner product, its spectrum lies in $[-1,1]$ (unital channels are contractions for the Hilbert--Schmidt norm), and
$\Theta_+$ is self-adjoint and positive semidefinite, as Proposition~\ref{prop:autonomous} requires of an autonomous limit.

\begin{proposition}[An expensive autonomous limit]\label{prop:lazy}
Let $\Phi$, $V$ and $\Theta$ be as in Remark~\ref{rem:doubling}, and $\kappa=6+\log107$.
\begin{enumerate}[label=(\alph*),itemsep=1pt]
\item \emph{Autonomy.} Let $Z$ be the coordinate of $L^\infty[-1,1]$ with the normalized Lebesgue measure, put
$\mathcal B=L(\Gamma\times\Z_3)\,\overline\otimes\,L^\infty[-1,1]$ with the product trace, and
$W=\exp\bigl(i\tfrac\pi2V\otimes Z\bigr)=\cos(\tfrac\pi2Z)+i\sin(\tfrac\pi2Z)\,V\in M_{1746}\otimes\mathcal B$. The autonomous dynamics
$\alpha=\mathrm{Ad}\,W$ has $\Phi_t=\Theta_+$ for every $t\ge1$. More generally, $\exp(isV\otimes Z)$ gives the channel
$\Xi_s=\frac{1+\mathrm{sinc}\,2s}2\,\mathrm{id}+\frac{1-\mathrm{sinc}\,2s}2\,\Theta$, with $\mathrm{sinc}\,x=\sin x/x$, and
$\ddia(\Xi_s,\Theta_+)=|\sin2s|/(4s)$.
\item \emph{Memory against purity, lower bound.} There are $c,c_0>0$ and $n_0$ such that for $n\ge n_0$ and $\log n\le B\le c_0n$ every causal
service of $\Theta_+$ for $n$ rounds with error at most $1/16$ and purity $P+C\le B$ has, at any exchange,
\[
Q\ \ge\ c\,\frac n{B\,(\log n)^{2\kappa}} .
\]
In particular, for every $A_0\ge0$ there is $c'>0$ such that services with $P+C\le A_0\log n$ have $Q\ge c'n/(\log n)^{2\kappa+1}$ for all large $n$.
\item \emph{Upper bound.} There is $K$ such that for every $n\ge1$ and $0<B\le n$, $\Theta_+$ has an exact causal service with purity $P+C<B$,
$J=(n-1)Q$ and $Q\le K(n/B)\log(2n/B)$. Moreover $\Theta_+\in\cK_{1746}$.
\end{enumerate}
\end{proposition}

So $\Theta_+$ is a long-time limit, indeed an exact finite-time value, of an autonomous dynamics with a tracial bath; it is expensive in the
sense defined at the beginning of this section, and it trades memory and purity at the same fixed rate as $\Phi$ (Theorem~\ref{thm:main-channel}(a)). Nothing in the construction is local. Part (b) loses nothing against
Theorem~\ref{thm:main-channel}(a) because the round extracted from a service of $\Theta_+$ is converted into a round of $\Phi$ with a
\emph{linear} loss in leakage and in the commutator $\|[U,1\otimes\sqrt\rho\,]\|_F^2$, which replaces entropy production in the one-round theorem
(Lemmas~\ref{lem:commutator-form} and~\ref{lem:block-extraction}), on the same bath.

\begin{proof}[Proof of (a)]
Since $V^2=1$ and $Z$ commutes with $V$, $\exp(isV\otimes Z)=\cos(sZ)+i\sin(sZ)V$, so that $W^t=\exp(it\frac\pi2V\otimes Z)$ and
\[
W^t(X\otimes1)W^{-t}=\cos^2(\tfrac{t\pi}2Z)\,X+\sin^2(\tfrac{t\pi}2Z)\,V(X\otimes1)V+\tfrac i2\sin(t\pi Z)\bigl(V(X\otimes1)-(X\otimes1)V\bigr).
\]
The trace of $\mathcal B$ is the product of the trace of $L(\Gamma\times\Z_3)$ and the integral over $Z$. The last term integrates to $0$ because
$\sin(t\pi Z)$ is odd, $\int\cos^2(\frac{t\pi}2Z)=\frac12(1+\mathrm{sinc}\,t\pi)=\frac12$ for $t\ge1$, and $(\mathrm{id}\otimes\tau)[V(X\otimes1)V]=\Theta(X)$.
So $\Phi_t=\frac12(\mathrm{id}+\Theta)$. The same computation with $s$ in place of $\frac{t\pi}2$ gives $\Xi_s$. Finally
$\Xi_s-\Theta_+=\frac12\mathrm{sinc}(2s)(\mathrm{id}-\Theta)$, and $\|\mathrm{id}-\Theta\|_\diamond=2$: the input $|0\rangle\langle0|\otimes\sigma$ is mapped
by $\Theta$ into the block $|1\rangle\langle1|$. The automorphism $\mathrm{Ad}\,W$ is inner, hence trace preserving.
\end{proof}

The same channel is reached, for every $t\ge1$, by a Bernoulli-shift construction like that of Proposition~\ref{prop:nonlocal}: with
$\mathcal B=L(\Gamma\times\Z_3)\,\overline\otimes\,\overline\bigotimes_{k\in\Z}(M_2,\tau_2)$, $\beta$ the shift of the chain and $p_0$ the projection $|1\rangle\langle1|$ of the qubit at site $0$, the automorphism
$\mathrm{Ad}\bigl((1-p_0)+p_0V\bigr)\circ(\mathrm{id}\otimes\beta)$ raises $V$ to the parity of $t$ fair bits.

We now work with a general unital channel $\Phi$ on $M_d$, its doubling $\Theta(\begin{smallmatrix}x&y\\u&w\end{smallmatrix})=\mathrm{diag}(\Phi(w),\Phi^*(x))$
and $\Theta_+=\frac12(\mathrm{id}+\Theta)$ on $M_{2d}$. Let $\mathcal S_\Phi\subseteq M_d$ be the span of the Kraus operators of $\Phi$, so that the Choi
support of $\Phi$ is $\{|A\rangle\!\rangle:A\in\mathcal S_\Phi\}$, and let $\Pi_\Phi$ also denote the Hilbert--Schmidt projection of $M_d$ onto
$\mathcal S_\Phi$. Since $\Theta(X)=\sum_g(e_{01}\otimes A_g)X(e_{01}\otimes A_g)^*+(e_{10}\otimes A_g^*)X(e_{10}\otimes A_g^*)^*$ for Kraus operators
$A_g$ of $\Phi$, the Choi support of $\Theta_+$ corresponds to the span of the three pairwise orthogonal subspaces
\begin{equation}\label{eq:lazy-support}
\mathbb C\,1_{2d},\qquad e_{01}\otimes\mathcal S_\Phi,\qquad e_{10}\otimes\mathcal S_\Phi^*
\end{equation}
of $M_{2d}=M_2\otimes M_d$. Throughout, $\|\cdot\|_F$ is the Hilbert--Schmidt norm and $\|\cdot\|$ the operator norm, as in
Section~\ref{sec:weighted-proof}; an operator $X$ on the bath $\mathbb C^M$ is identified with $1\otimes X$ when it acts on $\mathbb C^d\otimes\mathbb C^M$,
so that $\|1\otimes X\|_F=\sqrt d\,\|X\|_F$.

\begin{lemma}[Commutator form of the one-round theorem]\label{lem:commutator-form}
Lemma~\ref{lem:weighted-extraction}, Theorem~\ref{thm:weighted} and Corollary~\ref{cor:amplification} remain true if, in the definition
$\nu=\ell+\Delta S$, the entropy production $\Delta S$ of the round $\Psi(X)=\mathrm{Tr}_M\,U(X\otimes\rho)U^*$ is replaced by any $\tilde s\ge0$ with
$\|[U,1\otimes\sqrt\rho\,]\|_F^2\le4d\,\tilde s\ln2$, and if the hypothesis $\ddia(\Psi,\Phi_G)\le1/8$ is replaced by the weaker hypothesis
\begin{equation}\label{eq:coherence}
|\langle z|\Psi(|z\rangle\langle0|)|0\rangle|\le\tfrac34,
\end{equation}
where $0$ and $z$ are the anchor positions of the empty word and of $j$; in Lemma~\ref{lem:weighted-extraction}(iii) the bound then reads
$E(\phi(a)\phi(b))\ge0.66-e$. Conversely, every round can be rewritten, without changing its channel or its leakage, as a round
$(U',\rho)$ with $U'=(1\otimes W)U$ for which $\tilde s=\Delta S$ is admissible.
\end{lemma}

\begin{proof}
The rewriting is the normalization step of the proof of Lemma~\ref{lem:weighted-extraction}, after which~\eqref{eq:block-commutators} is
the inequality $\|[U',1\otimes\sqrt\rho\,]\|_F^2\le4d\,\Delta S\ln2$; its proof uses only that $U'$ is a unitary on $\mathbb C^d\otimes\mathbb C^M$. In
the proof of Lemma~\ref{lem:weighted-extraction}, $\Delta S$ enters only through~\eqref{eq:block-commutators}: afterwards only the bounds
$\epsilon_S\le2\sqrt{\ln2}\,\eta$, $\epsilon_\ell+\epsilon_S\le\sqrt{1+4\ln2}\,\eta$ and $2\epsilon_\ell^2+2\epsilon_S^2\le8\ln2\,\eta^2$ with $\eta=\sqrt{d\nu}$ are used,
together with the unitarity of $U$, its leakage and $\ddia(\Psi,\Phi_G)$. These bounds hold with $\epsilon_S^2\le4d\tilde s\ln2$ and $\nu=\ell+\tilde s$. The diamond hypothesis enters only in the displacement step, through
$|\Tr(\rho B_\varnothing^*B_j)|=|\langle z|\Psi(|z\rangle\langle0|)|0\rangle|\le2u\le\frac14$. Under~\eqref{eq:coherence} instead,
$E(\phi(j))^2\ge2-\frac32-4\epsilon_\ell\ge\frac12-4\eta$, so $E(\phi(a)\phi(b))\ge\sqrt{1/2-4\eta}-e\ge0.66-e$ if $\eta\le\frac1{64}$, while $e>1$ otherwise. In
the proof of Theorem~\ref{thm:weighted} this gives $E(Z_i)\ge0.66-\delta_L$ and $\Delta_i\ge0.66-30\delta_L\ge0.66-b_*>\frac12$ in place of
$\Delta_i\ge1-b_*$, and $\Delta_i\ge\frac12$ is all that Corollary~\ref{cor:sector} needs. Theorem~\ref{thm:weighted} and Corollary~\ref{cor:amplification} use
the round only through Lemma~\ref{lem:weighted-extraction}.
\end{proof}

\begin{lemma}[Weighted intertwining]\label{lem:twisted}
Let $B$ be an operator on $\mathbb C^d\otimes\mathbb C^M$ with $\|B\|\le1$, let $R,R'$ be self-adjoint operators on $\mathbb C^M$, $\xi$ an operator on $\mathbb C^M$, and
$\Sigma=BR-R'B$.
\begin{enumerate}[label=(\roman*),itemsep=1pt]
\item If $h$ is Lipschitz with constant $L_h$ on an interval containing the spectrum of $R$, then $\|[h(R),\xi]\|_F\le L_h\,\|[R,\xi]\|_F$.
\item If $g(x)=\int\hat g(t)e^{itx}dt$ with $M_1=\int|t\,\hat g(t)|\,dt$ and $M_2=\int t^2|\hat g(t)|\,dt$ finite, then
\[
\bigl\|\bigl(Bg(R)-g(R')B\bigr)\xi\bigr\|_F\ \le\ M_1\,\|\Sigma\xi\|_F+\tfrac12M_2\,\|\Sigma\|\,\sqrt d\,\|[R,\xi]\|_F .
\]
\end{enumerate}
\end{lemma}

\begin{proof}
(i) In an eigenbasis of $R$ with eigenvalues $r_k$, the entries of $[h(R),\xi]$ and $[R,\xi]$ are $(h(r_k)-h(r_l))\xi_{kl}$ and $(r_k-r_l)\xi_{kl}$.
(ii) Differentiating $\theta\mapsto e^{i\theta tR'}Be^{i(1-\theta)tR}$ gives
$Be^{itR}-e^{itR'}B=it\int_0^1e^{i\theta tR'}\,\Sigma\,e^{i(1-\theta)tR}\,d\theta$. Write $e^{i(1-\theta)tR}\xi=\xi e^{i(1-\theta)tR}+[e^{i(1-\theta)tR},\xi]$ and use
(i) with $h(x)=e^{i(1-\theta)tx}$; unitaries do not change Hilbert--Schmidt norms, so
$\|(Be^{itR}-e^{itR'}B)\xi\|_F\le|t|\,\|\Sigma\xi\|_F+\frac{t^2}2\|\Sigma\|\sqrt d\,\|[R,\xi]\|_F$. Integrate against $|\hat g(t)|\,dt$.
\end{proof}

We fix the cutoff once and for all. Let $S:\mathbb R\to[0,1]$ be smooth with $S=0$ on $(-\infty,0]$ and $S=1$ on $[1,\infty)$, let $0<s_0\le\frac14$, and put
\[
g(x)=\sin\bigl(\tfrac\pi2S(2x/s_0-1)\bigr)\,S(3-2x),\qquad g_2(x)=\cos\bigl(\tfrac\pi2S(2x/s_0-1)\bigr),\qquad f(x)=\sqrt x\,g(x).
\]
Then $g$ is smooth with compact support, so $M_g:=\int(|t|+t^2)|\hat g(t)|\,dt<\infty$; on $[0,1]$ we have $0\le g\le1$, $g=0$ on $[0,s_0/2]$,
$g=1$ and $g_2=0$ on $[s_0,1]$, and $g^2+g_2^2=1$. Let $L_g$ be a common Lipschitz constant of $g$ and $g_2$ on $[0,1]$, put $\sigma=\sqrt{2/s_0}$ and
\[
K_a=11M_g+1+1.5L_g+2\sigma(4+1.5L_g).
\]
These depend only on $s_0$.

\begin{lemma}[Block extraction]\label{lem:block-extraction}
Let $\Phi$ be a unital channel on $M_d$, and let $\Psi(X)=\mathrm{Tr}_M\,U(X\otimes\rho)U^*$ be a round of $\Theta_+$, with $U$ unitary on
$\mathbb C^2\otimes\mathbb C^d\otimes\mathbb C^M$, $u=\ddia(\Psi,\Theta_+)$ and leakage $\ell$ relative to $\Theta_+$. Put
$\epsilon_\ell=\sqrt{2d\ell}$, $\epsilon_S=\|[U,1\otimes\sqrt\rho\,]\|_F$, $\gamma=\epsilon_\ell+\epsilon_S$ and $\mathcal E=2u+s_0+4K_a\sqrt d\,\gamma$. If $\mathcal E<\frac12$,
there are a unitary $U'$ on $\mathbb C^d\otimes\mathbb C^M$ and a state $\rho'$ on $\mathbb C^M$ such that the round
$\Psi'(X)=\mathrm{Tr}_M\,U'(X\otimes\rho')U'^*$ of $\Phi$, with leakage $\ell_\Phi(\Psi')$ relative to $\Phi$, satisfies
\[
\ddia(\Psi',\Phi)\le\frac{\mathcal E}{\frac12-\mathcal E},\qquad \ell_\Phi(\Psi')\le\frac{4K_a^2\gamma^2}{d(\frac12-\mathcal E)},\qquad
\bigl\|[U',1\otimes\sqrt{\rho'}\,]\bigr\|_F^2\le\frac{8K_a^2\gamma^2}{\frac12-\mathcal E},
\]
and $|\langle x|\Psi'(|x\rangle\langle y|)|y\rangle|\le \mathcal E/(\frac12-\mathcal E)$ for all unit vectors $x,y$ with $\langle x|\Phi(|x\rangle\langle y|)|y\rangle=0$.
\end{lemma}

\begin{proof}
Put $\xi=\sqrt\rho$ and write $U=\sum_{a,b\in\{0,1\}}e_{ab}\otimes U_{ab}$ with $D_0=U_{00}$, $B=U_{01}$, $D_1=U_{11}$, operators on $\mathbb C^d\otimes\mathbb C^M$.

\emph{Residuals.} As in the proof of Lemma~\ref{lem:weighted-extraction}, $\mathcal J(\Psi)$ is the marginal of the unit vector
$(2d)^{-1/2}U\xi\in M_{2d}\otimes M_M$, so $\ell=(2d)^{-1}\|(\Pi^\perp\otimes\mathrm{id})(U\xi)\|_F^2$, with $\Pi$ the projection of $M_{2d}$ onto the span
of~\eqref{eq:lazy-support}. Put $\alpha=(2d)^{-1}(\Tr\otimes\mathrm{id})(U)$, the average of the $2d$ diagonal entries of $U$ in $M_M$, and
$E_a=D_a-\alpha$, $F=(\Pi_\Phi^\perp\otimes\mathrm{id})(B)$. By~\eqref{eq:lazy-support}, $(\Pi^\perp\otimes\mathrm{id})(U)$ has the diagonal blocks $E_0,E_1$ and
the upper right block $F$; since $\Pi^\perp\otimes\mathrm{id}$ commutes with right multiplication by $\xi$, the numbers $e_a=\|E_a\xi\|_F$ and
$e_F=\|F\xi\|_F$ are at most $\epsilon_\ell$. Also $\|\alpha\|\le1$ and $\|E_a\|\le2$.

\emph{Commutators.} Put $k_X=\|[X,\xi]\|_F$ for $X\in\{D_0,B,D_1\}$; these are at most $\epsilon_S$. The diagonal entries of $U$ are entries of $D_0$ and
$D_1$, so by Cauchy--Schwarz $k_\alpha:=\sqrt d\,\|[\alpha,\xi]\|_F\le(k_{D_0}^2+k_{D_1}^2)^{1/2}/\sqrt2\le\epsilon_S/\sqrt2$. With $R=1-\alpha^*\alpha$ and
$R'=1-\alpha\alpha^*$, both between $0$ and $1$, $k_R:=\sqrt d\,\|[R,\xi]\|_F\le2k_\alpha\le\sqrt2\,\epsilon_S$, and likewise $k_{R'}\le\sqrt2\,\epsilon_S$.

\emph{Defects.} Unitarity gives $B^*B=1-D_1^*D_1$ and $BB^*=1-D_0D_0^*$, hence
\begin{gather*}
\Delta:=R-B^*B=\alpha^*E_1+E_1^*D_1,\qquad\Delta':=R'-BB^*=\alpha E_0^*+E_0D_0^*,\\
\Sigma:=BR-R'B=B\Delta-\Delta'B,
\end{gather*}
with $\|\Delta\|,\|\Delta'\|\le1$ and $\|\Sigma\|\le2$. From $\xi E_1=E_1\xi-[D_1,\xi]+[\alpha,\xi]$ and $D_1\xi=\xi D_1+[D_1,\xi]$,
$\|\Delta\xi\|_F\le e_1+\|\xi E_1\|_F+2k_{D_1}\le2e_1+3k_{D_1}+k_\alpha\le4\gamma$; symmetrically $\|\Delta'\xi\|_F\le2e_0+3k_{D_0}+k_\alpha\le4\gamma$;
and $\Sigma\xi=B\Delta\xi-\Delta'\xi B-\Delta'[B,\xi]$ gives $\|\Sigma\xi\|_F\le9\gamma$.

\emph{Two factorizations.} Let $P$ and $P'$ be the spectral projections of $R$ and $R'$ for $(s_0/2,\infty)$. Since $g$ vanishes on $[0,s_0/2]$,
$g(R)=R^{-1/2}Pf(R)$ with $\|R^{-1/2}P\|\le\sigma$, and likewise for $R'$. The operators $B_1=BR^{-1/2}P$ and $B_1'=R'^{-1/2}P'B$ satisfy
\[
B_1^*B_1=P-\tilde\Delta,\quad \tilde\Delta=PR^{-1/2}\Delta R^{-1/2}P,\qquad B_1'B_1'^*=P'-\tilde\Delta',\quad\tilde\Delta'=P'R'^{-1/2}\Delta'R'^{-1/2}P'.
\]
Take polar decompositions $B_1=V_1|B_1|$ and $B_1'=|B_1'^*|V_1'$ with $V_1,V_1'$ unitary. On the range of $1\otimes P$ we have
$|B_1|=1+\varphi(\tilde\Delta)$ with $\varphi(x)=\sqrt{1-x}-1$ and $\tilde\Delta\le1$, and $\varphi(x)^2\le x^2$; so $\|(|B_1|-P)Y\|_F\le\|\tilde\Delta Y\|_F$
for every $Y$ with range in that subspace, and similarly for $B_1'$. Put $Z_0=Bg(R)\xi=B_1f(R)\xi$. Then
$Z_0-V_1f(R)\xi=V_1(|B_1|-P)f(R)\xi$, and $\tilde\Delta f(R)\xi=PR^{-1/2}\Delta g(R)\xi$, while
$\|\Delta g(R)\xi\|_F\le\|\Delta\xi\|_F+\sqrt d\,\|[g(R),\xi]\|_F\le\|\Delta\xi\|_F+L_g\,k_R$ by Lemma~\ref{lem:twisted}(i). Hence
\[
a_1:=\|Z_0-V_1f(R)\xi\|_F\le\sigma(4+1.5L_g)\gamma .
\]
On the other side,
$Z_0=(Bg(R)-g(R')B)\xi+g(R')[B,\xi]+[g(R'),\xi]B+\xi f(R')B_1'$ and $\xi f(R')B_1'=\xi f(R')V_1'+\xi f(R')(|B_1'^*|-P')V_1'$, where
$\|\xi f(R')(|B_1'^*|-P')\|_F\le\|\tilde\Delta'f(R')\xi\|_F\le\sigma(\|\Delta'\xi\|_F+L_g\,k_{R'})$. By Lemma~\ref{lem:twisted}(ii),
$\|(Bg(R)-g(R')B)\xi\|_F\le M_g(\|\Sigma\xi\|_F+k_R)\le11M_g\gamma$. Hence
\[
a_2:=\|Z_0-\xi f(R')V_1'\|_F\le11M_g\gamma+(1+1.5L_g)\gamma+\sigma(4+1.5L_g)\gamma,\qquad a_1+a_2\le K_a\gamma .
\]

\emph{Polar decomposition of $Z_0$.} Put $\zeta_1=|f(R)\xi|=(\xi f(R)^2\xi)^{1/2}$ and $\zeta_2=|(\xi f(R'))^*|=(\xi f(R')^2\xi)^{1/2}$. Then $|V_1f(R)\xi|=\zeta_1$ and $|(\xi f(R')V_1')^*|=\zeta_2$, as operators $1\otimes\zeta_i$. By the
Araki--Yamagami inequality $\||X|-|Y|\|_F\le\sqrt2\,\|X-Y\|_F$~\cite{ArakiYamagami1981}, applied to $Z_0$ and to $Z_0^*$,
$\||Z_0|-\zeta_1\|_F\le\sqrt2\,a_1$ and $\||Z_0^*|-\zeta_2\|_F\le\sqrt2\,a_2$. Let $Z_0=O|Z_0|=|Z_0^*|O$ with $O$ unitary. Then
\[
\|O\zeta_1-\zeta_2O\|_F\le\|O(\zeta_1-|Z_0|)\|_F+\|(|Z_0^*|-\zeta_2)O\|_F\le\sqrt2(a_1+a_2)=:\delta .
\]

\emph{The round.} The Hermitian operators $O(1\otimes\zeta_1)O^*$ and $1\otimes\zeta_2$ are within $\delta$ in $\|\cdot\|_F$, and their spectra are those of $\zeta_1$
and $\zeta_2$, each repeated $d$ times. By the Hoffman--Wielandt inequality~\cite{HoffmanWielandt1953}, the decreasing eigenvalue lists
$(x_k)$ of $\zeta_1$ and $(y_k)$ of $\zeta_2$ satisfy $d\sum_k(x_k-y_k)^2\le\delta^2$; so the unitary $W_h$ of $\mathbb C^M$ that maps an eigenbasis of
$\zeta_1$ to one of $\zeta_2$ in the same order has $\|1\otimes(W_h\zeta_1W_h^*-\zeta_2)\|_F\le\delta$. Put $U'=(1\otimes W_h^*)O$, $c_1=\Tr\zeta_1^2=\Tr\rho f(R)^2$
and $\rho'=\zeta_1^2/c_1$. Then $\|[U',1\otimes\zeta_1]\|_F\le\|O\zeta_1-\zeta_2O\|_F+\|(\zeta_2-W_h\zeta_1W_h^*)O\|_F\le2\delta$, which is the third bound once we
know $c_1\ge\frac12-\mathcal E$. The channel of $(U',\rho')$ is $\Psi'=\Phi_{O\zeta_1}/c_1$, where $\Phi_X(w)=\mathrm{Tr}_M[X(w\otimes1)X^*]$. Its leakage is
$(dc_1)^{-1}\|(\Pi_\Phi^\perp\otimes\mathrm{id})(O\zeta_1)\|_F^2$, and
\[
\|(\Pi_\Phi^\perp\otimes\mathrm{id})(O\zeta_1)\|_F\le\|(\Pi_\Phi^\perp\otimes\mathrm{id})(Z_0)\|_F+\sqrt2\,a_1\le e_F+L_g\,k_R+\sqrt2\,a_1\le2K_a\gamma,
\]
since $(\Pi_\Phi^\perp\otimes\mathrm{id})(Bg(R)\xi)=F\xi\,g(R)+(\Pi_\Phi^\perp\otimes\mathrm{id})(B[g(R),\xi])$. This is the second bound.

\emph{Closeness.} For a completely positive map $\Phi_X$ as above, $\|\Phi_X\|_\diamond=\|\mathrm{Tr}_M(X^*X)\|$, and
$\|\Phi_X-\Phi_Y\|_\diamond\le\|X-Y\|_F\bigl(\|\mathrm{Tr}_M X^*X\|^{1/2}+\|\mathrm{Tr}_M Y^*Y\|^{1/2}\bigr)$, where $\mathrm{Tr}_M$ traces out the input bath.
Let $Y=B\xi$. The upper left block of $\Psi(e_{11}\otimes w)$ is $\Phi_Y(w)$ and that of $\Theta_+(e_{11}\otimes w)$ is $\frac12\Phi(w)$; compressing to
the block, $\|\Phi_Y-\frac12\Phi\|_\diamond\le2u$. Since $g^2+g_2^2=1$ on $[0,1]$, $\Phi_Y-\Phi_{Z_0}-\Phi_{Bg_2(R)\xi}$ is
$w\mapsto\mathrm{Tr}_M[B(w\otimes X)B^*]$ with $X=\sum_{h\in\{g,g_2\}}h(R)[h(R),\rho]$, and $\|X\|_1\le4L_g\|[R,\xi]\|_F$, so its diamond norm is at most
$4L_g\,k_R/\sqrt d$ (split $X$ into positive and negative parts and use $B^*B\le1$). Next,
$\|\Phi_{Bg_2(R)\xi}\|_\diamond\le\Tr\rho\,g_2(R)Rg_2(R)+\|\xi g_2(R)\Delta g_2(R)\xi\|_1\le s_0+\sqrt d(\|\Delta\xi\|_F+L_g\,k_R)$, because $xg_2(x)^2\le s_0$ on
$[0,1]$. Finally $\mathrm{Tr}_M\zeta_1^2=c_1\le1$ and $\mathrm{Tr}_M Z_0^*Z_0\le1$ give $\|\Phi_{O\zeta_1}-\Phi_{Z_0}\|_\diamond\le2\sqrt2\,a_1$. Adding up,
\[
\bigl\|\Phi_{O\zeta_1}-\tfrac12\Phi\bigr\|_\diamond\le2u+s_0+2\sqrt2\,a_1+\sqrt d(4+1.5L_g)\gamma+6L_g\gamma/\sqrt d\le \mathcal E .
\]
Evaluating at $\tau_d$ gives $|c_1-\frac12|\le \mathcal E$, so $c_1\ge\frac12-\mathcal E$, and $\|\Psi'-\Phi\|_\diamond\le \mathcal E/c_1+|1/(2c_1)-1|\le2\mathcal E/c_1$. If
$\langle x|\Phi(|x\rangle\langle y|)|y\rangle=0$, then $|\langle x|\Phi_{O\zeta_1}(|x\rangle\langle y|)|y\rangle|\le\|\Phi_{O\zeta_1}-\frac12\Phi\|_\diamond\le \mathcal E$, and dividing by
$c_1$ gives the last bound.
\end{proof}

\begin{proof}[Proof of (b)]
Let a service of $\Theta_+$ for $n$ rounds have error $\epsilon\le\frac1{16}$ and $P+C\le B$ with $\log n\le B$. Theorem~\ref{thm:extraction},
applied to $\Theta_+$ on $M_{1746}$, gives a round with bath $\mathbb C^M$, $M=2^Q$, $u=\ddia(\Psi_t,\Theta_+)\le\frac1{15}$ and
$\nu_t=\ell_t+\Delta S_t\le t_n\le\frac{16}{15}\cdot\frac{B+1+\log n}n\le\frac{3.2\,B}n$, as in the proof of Theorem~\ref{thm:main-channel}. The normalization
step of the proof of Lemma~\ref{lem:weighted-extraction}, which uses only that $U_t$ is a unitary on $\mathbb C^{2d}\otimes\mathbb C^M$, replaces $U_t$ by
$(1\otimes W)U_t$ without changing $\Psi_t$ or its leakage, after which $\epsilon_S^2=\|[U_t,1\otimes\sqrt{\rho_t}\,]\|_F^2\le4(2d)\Delta S_t\ln2$; so
$\gamma^2\le2(2d\ell_t+8d\ln2\,\Delta S_t)\le12d\,\nu_t$. Apply Lemma~\ref{lem:block-extraction} with $s_0=\frac1{100}$ and $d=873$:
$\mathcal E\le\frac2{15}+\frac1{100}+14K_ad\sqrt{\nu_t}\le\frac3{14}$ as soon as $14K_ad\sqrt{3.2B/n}\le0.07$. Then $c_1\ge\frac12-\mathcal E\ge\frac27$, and the anchors $0$
and $z=z(j)$ satisfy $\langle z|\Phi(|z\rangle\langle0|)|0\rangle=\tau(\lambda([j]))=0$, so the round $\Psi'$ satisfies~\eqref{eq:coherence}. By
Lemma~\ref{lem:commutator-form} it may be used in Corollary~\ref{cor:amplification} with
\[
\nu'=\ell_\Phi(\Psi')+\frac{\|[U',1\otimes\sqrt{\rho'}]\|_F^2}{4d\ln2}\le\Bigl(4+\frac2{\ln2}\Bigr)\frac{K_a^2\gamma^2}{d\,c_1}\le290K_a^2\,\nu_t\le928K_a^2\,\frac Bn .
\]
Let $C_R$ and $y_0$ be the constants of Corollary~\ref{cor:amplification}(a) for $R_\Phi$ and the transport of Section~\ref{sec:lamp-channel}, and put
$y_n=n/(928K_a^2C_RB)$, so that $C_R\nu'y_n\le1$ and $y_n\le n$. With
$c_0=\min\{1/(928K_a^2C_Ry_0),\ (0.07/(14K_ad))^2/3.2\}$ and $B\le c_0n$, the corollary gives
\[
Q=\log M\ \ge\ S(\rho')\ \ge\ \frac{y_n}{32(\log y_n)^{2\kappa}}\ \ge\ \frac1{29696\,K_a^2C_R}\cdot\frac n{B(\log n)^{2\kappa}} .
\]
For the second statement take $B=\max(1,A_0)\log n$.
\end{proof}

\begin{proof}[Proof of (c)]
Let $E$, $K_G$ and a $(1-r^{-1})$-expansive $r^{-1}$-morphism $f:E\to\Sym(\Omega)$ be as in Proposition~\ref{prop:sofic-upper}, with its unitary
$U_f=\sum_ac_aE_a\otimes P_{\pi_a}$ and its set $\Omega_0$ of good points, at which $\mathrm{Tr}_\Omega U_f(w\otimes|\omega\rangle\langle\omega|)U_f^*=\Phi(w)$. Let
$\Omega_1$ be the set of $\omega$ with $(f([x])f([y]))^{-1}(\omega)=f([z])^{-1}(\omega)$ for every triangle $(x,y,z)$ and $f(g)^{-1}(\omega)\ne f(h)^{-1}(\omega)$ for
distinct labels $g,h$. As $U_f^*(\psi\otimes|\omega\rangle)=\sum_a\overline{c_a}E_a^*\psi\otimes|\pi_a^{-1}(\omega)\rangle$, the computation of
Proposition~\ref{prop:sofic-upper} gives $\mathrm{Tr}_\Omega U_f^*(x\otimes|\omega\rangle\langle\omega|)U_f=\Phi^*(x)$ for $\omega\in\Omega_1$. A point outside $\Omega_1$
is the image under the bijection $f([z])$ of a point where $f([x])f([y])$ and $f([z])$ differ, or under $f(g)$ of a point where $f(g)$ and $f(h)$
agree; so $|\Omega\setminus\Omega_1|\le K_GM/r$, and $\Omega_*=\Omega_0\cap\Omega_1$ has $|\Omega_*|\ge(1-2K_G/r)M$.

Let the bath be a coin qubit, a phase qubit carrying $z=\mathrm{diag}(1,i)$, and $\mathbb C^\Omega$, and put
$W_f=|0\rangle\langle0|\otimes1+|1\rangle\langle1|\otimes V_f$ with $V_f=\bigl(\begin{smallmatrix}0&U_f\otimes z\\U_f^*\otimes z^*&0\end{smallmatrix}\bigr)$. For every
state $\beta$ that is a mixture of the $|\omega\rangle\langle\omega|$, $\omega\in\Omega_*$, the bath state $\tau_2\otimes\tau_2\otimes\beta$ makes $W_f$ apply
exactly $\Theta_+$: the coin splits the channel into $\frac12\mathrm{id}$ and $\frac12$ the channel of $V_f$, the phase qubit removes the off-diagonal
blocks since $\tau(z^2)=0$, and the diagonal blocks are $\Phi(w)$ and $\Phi^*(x)$ at the points of $\Omega_*$. Now repeat the exact service of
Proposition~\ref{prop:sofic-upper}(c) with $K_G$ replaced by $2K_G$ and two more qubits, which are imported maximally mixed: with
$\eta=B/(4n)$ and $r=2K_G/\eta$ the imported register is maximally mixed on $\mathrm{span}\{|\omega\rangle:\omega\in\Omega_*\}\otimes\mathbb C^{k_0}$ (and on the
two qubits), its deficit is below $B/n$, every round applies exactly $\Theta_+$, and $Q=\lceil\log\sigma_E(8K_Gn/B)+\log(4n/B)\rceil+2$. The bound
$\log\sigma_E(R)\le C_ER\log R$ of Theorem~\ref{thm:main-group}(b) gives $Q\le K(n/B)\log(2n/B)$. With $\beta$ maximally mixed on all of $\Omega$,
the same unitary gives a finite tracial factorization of $(1-\vartheta)\Theta_++\vartheta\,\Upsilon$ for some channel $\Upsilon$ and some $\vartheta\le2K_G/r$, for every $r$; hence
$\Theta_+\in\cK_{1746}$.
\end{proof}

\begin{remark}[Exact finite factorizations of $\Theta_+$]\label{rem:lazy-exact}
$\Theta_+$ is an exact autonomous value but not an exact finite factor. Indeed, if $\Theta_+(X)=(\mathrm{id}\otimes\tau_m)[W(X\otimes1)W^*]$ for a unitary
$W$ on $\mathbb C^{2d}\otimes\mathbb C^m$, the proof of Lemma~\ref{lem:block-extraction} with $\rho=\tau_m$ has $\ell=0$ and $\epsilon_S=0$: the diagonal blocks of
$W$ are both $1\otimes\alpha$, its upper right block $B$ lies in $\mathcal S_\Phi\otimes M_m$, $B^*B=1\otimes R$ and $BB^*=1\otimes R'$. Then
$B(1\otimes R)=(1\otimes R')B$, so for every eigenvalue $\mu>0$ of $R$ the operator $\mu^{-1/2}B$ maps $\mathbb C^d\otimes\ker(R-\mu)$ unitarily onto
$\mathbb C^d\otimes\ker(R'-\mu)$, and the upper left block of $\Theta_+(e_{11}\otimes w)$ writes $\Phi$ as the convex combination
$\sum_\mu\frac{2\mu m_\mu}m\,\Phi_\mu$ of exact factors with baths of dimension $m_\mu=\dim\ker(R-\mu)\le m$. For the channel of
Theorem~\ref{thm:main-channel} this is impossible: such a combination is a single round with $\ell=\Delta S=0$ (a direct sum of flat baths), and
Theorem~\ref{thm:weighted} with $\nu=0$ would give it entropy at least $N/16$ for every $N$.
\end{remark}

\begin{remark}[A lower exchange]\label{rem:lazy-exchange}
At an anchor $z$ of the gadget, $\Phi(|z\rangle\langle z|)=|z\rangle\langle z|$, so $\Theta_+$ maps $|1\rangle\langle1|\otimes|z\rangle\langle z|$ to the flat state
$\frac12(|0\rangle\langle0|+|1\rangle\langle1|)\otimes|z\rangle\langle z|$ of rank two. Proposition~\ref{prop:purity-budget} with $r=2$ therefore gives $P+C<3Q+6$ for
services of $\Theta_+$ with error at most $1/16$ and exchange $J\le n/2$, and Proposition~\ref{prop:lazy}(b), used as in the proof of
Theorem~\ref{thm:main-channel}(b) with $B=\max(3Q+6,\log n)$ (if $B>c_0n$, then already $Q\ge(c_0n-6)/3$), gives $Q\ge c\sqrt n/(\log n)^\kappa$ for such services, whatever their purity.
\end{remark}

So autonomous dynamics with a tracial bath can produce expensive channels. The reduced dynamics of the probe matters: if the channels
$\Phi_t$ themselves form a semigroup of unital channels, they are mixed unitary for all large $t$~\cite{BhatDevendra2025}, hence cheap, whereas
$\Theta_+^2\ne\Theta_+$. Locality constrains how $\Theta_+$ can be approached: Proposition~\ref{prop:light-cone} and
Proposition~\ref{prop:lazy}(b), combined as in the proof of Corollary~\ref{cor:preparation}, show that a local reservoir with $d=1746$ that
produces a channel within $\delta\le\delta_0$ of $\Theta_+$ at time $T$ has $T\ge c\,\delta^{-1/\nu}(\log(1/\delta))^{-(2\kappa+1)/\nu}$, for constants
$c,\delta_0>0$ that depend only on the parameters of the reservoir. So no local reservoir produces $\Theta_+$ exactly at a finite time, or approaches
it exponentially fast. Whether a local dynamics can produce an expensive channel at all is open (Section~\ref{sec:open}).
